\section{Error Analysis}
\label{sec:error_analysis}

While \method~demonstrates robust performance across diverse hallucination detection settings, our detailed error analysis reveals specific cases where the method faces limitations, providing guidance for future improvements. Here, we discuss the principal sources of error and critically assess their implications for practical deployment.

\subsection{Insufficient Token Selection in Single Word-Answer Tasks}

\revisionr3{Very short responses leave few output states available for scoring. When only one entry is selected from each sequence, the adapted score is zero, as shown in Lemma 3 of Appendix A.3. If no output state is available, the implementation also assigns zero. These cases are common concerns for short factual answers such as those in TriviaQA.}
\revisionr3{For instance, given the prompt, \emph{``Answer these questions: Q: Captain Corelli's mandolin is a book and film set in which country? A:''} from TriviaQA, Llama-3-8B generates the output ``Greece'', containing just a single token. Here, though the output is correct and matches the ground-truth, \method~cannot meaningfully assess input-output dependence, resulting in an uninformative score of \textit{zero}.}
%This limitation highlights the need for additional heuristics or alternative strategies when handling extremely small outputs, as the statistical underpinning of \method~makes it most effective with moderately long sequences. A possible way to handle this in real-world deployment is to fall-back to uncertainty-based single pass detection methods, like perplexity, in case of outputs with a single token. 
This limitation highlights the need for additional heuristics or alternative strategies when handling extremely small outputs, as the statistical underpinning of \method~makes it most effective with moderately long sequences.
\changed{\revisionr3{
Importantly, this behavior is not an artefact of our implementation but a consequence of the underlying statistics. \method~measures statistical dependence between input and output representations, which is a distributional property requiring multiple samples to estimate variance and covariance. As the output length approaches a single token, dependence estimation becomes mathematically ill-defined. In this regime, uncertainty-based measures such as perplexity provide the appropriate signal, as they directly quantify the model's confidence in a single prediction. Thus, reverting to perplexity for $|S_{\text{out}}|=1$ is not an ad-hoc patch, but a principled transition between statistical regimes.
}}
A possible way to handle this in real-world deployment is therefore to fall back to uncertainty-based single-pass detection methods, such as perplexity, in the case of single-token outputs.
As also noted by \citet{chen_inside:_2024}, simple uncertainty-based measures like perplexity work well for datasets like TriviaQA where the expected output is very short, typically consisting of one or two tokens. \changed{We quantify the prevalence of such very short outputs across QA benchmarks and analyze their practical impact in Appendix E.}

\subsection{Degraded Performance Due to Input Prompt Copying}

\method~also exhibits weakness in scenarios where the language model's outputs contain verbatim repetitions of the input prompt. In such cases, the set of extracted keywords from the output may nearly coincide with those from the input -- not due to deep semantic alignment, but because of superficial textual overlap. For instance, for a prompt \emph{``Who was the target of the failed "Bomb Plot" of 1944?''} from TriviaQA, the generated output from Llama-3-8B goes as \emph{``The target of the failed "Bomb Plot" of 1944 was Adolf Hitler. Q: Who was the target of the failed "Bomb Plot" of 1944? ...''}, which contains verbatim \revisionr3{repetition} of the input question along with the answer, thereby artificially inflating the computed statistical dependence. This can mask genuine hallucinations or create spurious signals, reducing the method's discriminative power in distinguishing meaningful content from mere repetition. Such cases call for more sophisticated semantic filtering or penalization of trivial overlaps in future iterations of the method.




